\item Rifting \\
A simple model of rifting is illustrated in Figure~\ref{fig:rifting}.
\begin{figure}[hbt]
    \centering
    \includegraphics[]{figures/rifting.eps}
    \caption{Temperature}
    \label{fig:rifting}
\end{figure}
\begin{enumerate}
    \item Assuming the rifting takes place instantaneously, what are the
    initial and equilibrium geotherms (temperature as a function of depth)?

    {\bf Answer:} \\
    To write down the initial condition and equilibrium conditions we need
    to look for a function to describe our drawing.  Since both the initial
    equilibrium geotherms are lines we use the form $y = b + m x$ to
    describe our line, where $y$ is the temperature and $x$ is depth.
    Hence, we need to determine the intercept and slope for each geotherm.
    At $z = 0$, both geotherms are $T_0$, which gives us our intercept.  The
    slope is define by a change in temperature over a change in depth.  To
    define the slope, we choose convenient points at the beginning and end
    of our lines.  Hence our geotherms are given by,
    \begin{align}
        f(z) &=
        \begin{cases}
            T_0 + (T_m - T_0) \frac{z}{L(1 - \gamma)} 
            & 0 \leq z < L(1 - \gamma) \\
            T_m & L(1 - \gamma) \leq z \leq L
        \end{cases} \\
        T_E(z) & = T_0 + (T_m - T_0) \frac{z}{L}.
    \end{align}
    The initial condition has cases because the geotherm changes at the base
    of the rifted lithosphere.  

    \item As discussed in the practical and notes, the Fourier coefficients
    for this type of problem are given by
    \begin{equation}
        a_n = \frac{2}{L} 
        \int_0^L [f(z) - T_E(z)] \sin \left( \frac{n \pi z}{L} \right) \dd{z} .
    \end{equation}
    Using the geotherms determined above, what are the values of the
    Fourier coefficients?

    {\bf Answer:}\\
    To derive the Fourier coefficients, we need to input the result from the
    previous part into the equation above.  Because $f(z)$ has cases we will
    need to perform a separate integration for each case (i.e. for $[0,L(1 -
    \gamma))$ and $[L(1 - \gamma),L]$),
    \begin{align*}
        a_{n} &= \frac{2}{L} \Biggl\{ 
            \int_0^{L(1 - \gamma)}
            \Biggl[\biggl(T_0 + (T_m - T_0) \frac{z}{L(1 - \gamma)}\biggr)
            - \biggl(T_0 + (T_m - T_0) \frac{z}{L}\biggr) \Biggr]
            \sin \arg{z} \dd{z} \Biggr. \\
            &+ \Biggl. \int_{L(1 - \gamma)}^L
            \Biggl[T_m - \biggl(T_0 + (T_m - T_0) \frac{z}{L}\biggr) \Biggr]
            \sin \arg{z} \dd{z}
        \Biggr\} .
    \end{align*}
    Before performing integrating, we will simplify terms by pulling out a
    $(T_m - T_0)$ since it appears in every term,
    \begin{align*}
        a_{n} &= \frac{2}{L} (T_m - T_0) \Biggl\{ 
            \int_0^{L(1 - \gamma)}
            \biggl(\frac{z}{L(1 - \gamma)} - \frac{z}{L}\biggr)
            \sin \arg{z} \dd{z} \Biggr. \\
            & + \Biggl. \int_{L(1 - \gamma)}^L
            \biggl( 1 - \frac{z}{L} \biggr) \sin \arg{z} \dd{z} \Biggr\}
    \end{align*}
    A $\frac{z}{L}$ term appears within each set of integral limits.
    Because the function is continuous across the domain boundaries we can
    combine each part into a single integral across the combined domain,
    \begin{align*}
        a_{n} &= \frac{2}{L} (T_m - T_0) \Biggl\{ 
            \frac{1}{L(1 - \gamma)} \underbrace{\int_0^{L(1 - \gamma)}
            z \sin \arg{z} \dd{z}}_1 \Biggr.\\
            & + \Biggl. \underbrace{\int_{L(1 - \gamma)}^L \sin \arg{z} \dd{z}}_2
            - \frac{1}{L} \underbrace{\int_0^L z \sin \arg{z} \dd{z}}_3 \Biggr\}
    \end{align*}
    Now we can look up and solve each of the integrals individually.  The
    integrals fall into one of two forms which have the solutions:
    \begin{align*}
        \int \sin (az) \dd{z} &= -\frac{1}{a} \cos (az) \\
        \int z \sin (az) \dd{z} &= \frac{1}{a^2} \sin (az) - \frac{z}{a} \cos
        (az) .
    \end{align*}
    Choosing $a = \frac{n \pi}{L}$ we can see how our integrals relate to
    the general forms given above.
    Evaluating each of the integrals,
    \begin{align*}
        1: && \int_0^{L(1 - \gamma)}
          z \sin \biggl( \frac{n \pi z}{L} \biggr) \dd{z}
          &= \frac{L^2}{n^2 \pi^2} \sin \arg{z}
          - \frac{L}{n \pi} z \cos \arg{z}
          \Biggl. \Biggr|_0^{L(1 - \gamma)} \\
          && &= \frac{L^2}{n^2 \pi^2} \Biggl[ \sin \arg{L(1 - \gamma)}
          + \sin (0) \Biggr] \\
          && &- \frac{L}{n \pi} \Biggl[ L(1 - \gamma) \cos \arg{L(1 - \gamma)}
          + 0 \cos (0) \Biggr] \\
          && &= \frac{L^2}{n^2 \pi^2} \sin (n \pi (1 - \gamma))
          - \frac{L^2}{n \pi} (1 - \gamma) \cos (n \pi (1 - \gamma)) ;\\
        2: && \int_{L(1 - \gamma)}^L
          \sin \arg{z} \dd{z} &= - \frac{L}{n \pi} \cos \arg{z}
          \Biggl. \Biggr|_{L(1 - \gamma)}^{L} \\
          && &= - \frac{L}{n \pi} \Biggl[ \cos \arg{L} - \cos \arg{L(1 -
          \gamma)}\Biggr] \\
          && &= - \frac{L}{n \pi} [ \cos (n \pi) - \cos (n \pi (1 - \gamma)) ] \\
          && &= - \frac{L}{n \pi} (-1)^n + \frac{L}{n \pi} \cos (n \pi (1
          - \gamma)) ;\\
        3: && \int_0^L z \sin \arg{z} \dd{z}
          &= \frac{L^2}{n^2 \pi^2} \sin \arg{z} - \frac{L}{n \pi} z \cos \arg{z}
          \Biggl. \Biggr|_0^{L} \\
          && &= \frac{L^2}{n^2 \pi^2} \Biggl[ \sin \arg{L} + \sin (0) \Biggr] 
          - \frac{L}{n \pi} \Biggl[ L \cos \arg{L} + 0 \cos (0) \Biggr] \\
          && &= \frac{L^2}{n^2 \pi^2} \sin (n \pi) - \frac{L^2}{n \pi} \cos
          (n \pi) \\
          && &= - \frac{L^2}{n \pi} (-1)^n .
    \end{align*}
    The following trigonometric relationships are helpful when evaluating
    the integration limits:
    \begin{align*}
        \sin (n \pi) &= 0 & \cos (n \pi) &= (-1)^n .
    \end{align*}
    With the integrals evaluated, we can plug their solutions back in to
    $a_n$,
    \begin{align*}
        a_{n} &= \frac{2}{L} (T_m - T_0) \Biggl\{ \Biggr.
          \frac{1}{L(1 - \gamma)} \biggl[ \frac{L^2}{n^2 \pi^2}
          \sin (n \pi (1 - \gamma))
          - \frac{L^2}{n \pi} (1 - \gamma) \cos (n \pi (1 - \gamma)) \biggr] \\
          &+ \biggl[ - \frac{L}{n \pi} (-1)^n + \frac{L}{n \pi} \cos (n \pi (1
          - \gamma)) \biggr]
          - \frac{1}{L} \biggl[- \frac{L^2}{n \pi} (-1)^n \biggr]
          \Biggl. \Biggr\} \\
        &= 2 (T_m - T_0) \Biggl\{ \Biggr.
          \frac{1}{n^2 \pi^2 (1 - \gamma)} \sin (n \pi (1 - \gamma))
          - \frac{1}{n \pi} \cos (n \pi (1 - \gamma)) \\
          &- \frac{1}{n \pi} (-1)^n + \frac{1}{n \pi} \cos (n \pi (1 - \gamma))
          + \frac{1}{n \pi} (-1)^n
          \Biggl. \Biggr\} \\
        &= 2 (T_m - T_0) \frac{\sin (n \pi (1 - \gamma))}{n^2 \pi^2 (1 -
        \gamma)} .
    \end{align*}
    We could stop at this point since it is sufficiently simplified, or we
    could apply a trigonometric identity,
    \begin{equation*}
        \sin (\alpha - \beta) = \sin \alpha \cos \beta - \cos \alpha \sin
        \beta,
    \end{equation*}
    since the angle angle $n \pi (1 - \gamma)$ represents a difference.
    Choosing $\alpha = n \pi$ and $\beta = n \pi \gamma$, we find
    \begin{align*}
        a_n &= 2 (T_m - T_0) 
            \frac{\sin (n \pi) \cos (n \pi \gamma) 
            - \cos (n \pi) \sin(n \pi \gamma) }
            {n^2 \pi^2 (1 - \gamma)} \\
        &= 2 (T_m - T_0) 
            \frac{ - (-1)^n \sin(n \pi \gamma) }
            {n^2 \pi^2 (1 - \gamma)} \\
        &= 2 (T_m - T_0) 
            \frac{ (-1)^{n+1} \sin(n \pi \gamma) }
            {n^2 \pi^2 (1 - \gamma)}
    \end{align*}


    \item There are two timescales for this problem: one as cooling
    progresses from the base of the rifted lithosphere to the surface; and a
    second as cooling progresses from the base of the rifted lithosphere to
    the base of the equilibrium lithosphere.  What are these timescales?

    {\bf Answer:} \\
    To find the timescale for cooling, we take the argument of the time dependent term of the solution,
    \begin{equation*}
        T(z,t) = T_E(z) + \sum_{n=1}^\infty a_n \sin \left( \frac{n \pi z}{L}
        \right) \underbrace{\exp \left( - \frac{n^2 \pi^2}{L^2} \kappa t
        \right)}_{t \text{-dependence}} .
    \end{equation*}
    Just as we have done for the half-space cooling model, we can define our time scale using the argument of this term set to some arbitrary value (we will choose 1 for simplicity),
    \begin{equation*}
        \frac{n^2 \pi^2}{L^2} \kappa t = 1,
    \end{equation*}
    We will assume $n = 1$, that's the dominant term in the sum, which gives
    \begin{align*}
        \frac{\pi^2}{L^2} \kappa t &= 1 \\
        t &= \frac{L}{\pi^2 \kappa},
    \end{align*}
    To find our time scale for cooling from the base of the rifted
    lithosphere, we replace $L$ with the relevant distances: 0 to $L(1 -
    \gamma)$; and $L(1 - \gamma)$ to $L$,
    \begin{align*}
        \tau_{\text{upper}} &= \frac{L(1 - \gamma)}{\pi^2 \kappa} \\
        \tau_{\text{lower}} &= \frac{\gamma L}{\pi^2 \kappa}.
    \end{align*}

    \item What happens as $\gamma \to 1$?  How does the timescale compare
    with this compare with that estimate by plate cooling?  Show your work.

    {\bf Answer:}
    As $\gamma \to 1$, $\tau_{\text{upper}} \to 0$ whereas
    $\tau_{\text{lower}} \to \frac{L^2}{\pi^2 \kappa}$, which is the same as
    the plate cooling.
\end{enumerate}
